\documentclass[a4paper,UTF8]{article}
\usepackage[T1]{fontenc}
\usepackage{CJKutf8}
\usepackage{geometry}
\geometry{verbose,tmargin=3cm,bmargin=3cm,lmargin=2cm}
\usepackage{amsmath}
\usepackage{amsthm}
\usepackage{amssymb}

\makeatletter
%%%%%%%%%%%%%%% Textclass specific LaTeX commands.
% amsthdoc.pdf manual
\theoremstyle{plain}
\newtheorem{theorem}{Theorem}[section]
\newtheorem{lemma}{Lemma}
\newtheorem{corrolary}{Corrolary}

\theoremstyle{definition}
\newtheorem{definition}{Definition}[section]

%%%%%%%%%%%%%%% User specified LaTeX commands.
% 如果没有这一句命令，XeTeX会出错，原因参见
% http://bbs.ctex.org/viewthread.php?tid=60547
\DeclareRobustCommand\nobreakspace{\leavevmode\nobreak\ }
\usepackage{tkz-euclide}
%\usetkzobj{all}

\def\degree{^\\circ}
\def\GF{\mathrm{GF}}
\def\Re{\mathrm{Re}}
\def\Im{\mathrm{Im}}
\def\Res{\mathrm{Res}}
\def\diag{\,\mathrm{diag}\,}
\def\sumcyc{\sum\limits_{cyc}}
\def\prodcyc{\prod\limits_{cyc}}
\def\\i{\\infty}

%%%%%%%%%%%%%%%%%%%%%%%%%%%%%%%%%%%%%%%%%%
%%%%%%%%%%%%%%%%%%%%%%%%%%%%%%%%%%%%%%%%%%
%%%%%%%%%%%%%%%%%%%%%%%%%%%%%%%%%%%%%%%%%%

% 放在导言区
\newif\ifshowanswers
%\showanswersfalse % 默认不显示答案；改成 \showanswerstrue 则显示答案
\showanswerstrue % 默认不显示答案；改成 \showanswerstrue 则显示答案

% \blank[横线宽度]{答案}
% 不显示答案时：画固定宽度横线
% 显示答案时：在横线上居中显示答案
\newcommand{\blank}[2][2cm]{%
	\ifshowanswers
	\underline{\makebox[#1][c]{#2}}%
	\else
	\underline{\hspace*{#1}}%
	\fi
}

%%%%%%%%%%%%%%%%%%%%%%%%%%%%%%%%%%%%%%%%%%
%%%%%%%%%%%%%%%%%%%%%%%%%%%%%%%%%%%%%%%%%%
%%%%%%%%%%%%%%%%%%%%%%%%%%%%%%%%%%%%%%%%%%

\makeatother
\title{全国高中数学联赛模拟试题精选}

\begin{document}
\begin{CJK}{UTF8}{gbsn}
	
	\maketitle
	
	\section{模拟试题 1}
	\subsection{第一试}
	\subsubsection*{一、填空题（本题满分 64 分，每小题 8 分）}
	
	1．在金属丝制作的 $3 \times 4 \times 7$ 的长方体框架中放置着一个球，则该球的半径的最大值为\blank{$5/2$}．
	
	2．双曲线的左、右两焦点分别为 $F_1, F_2$ ，一条过 $F_2$ 的直线与双曲线的右支交于 $A, B$两点．若 $\triangle F_1 A B$ 是正三角形，则双曲线的离心率是 \blank{$\sqrt{3}$}．
	
	3．已知正实数 $a, b, c$ 满足 $a+b+c=1$ ，则 $a+\sqrt{b}+\sqrt[3]{c}$ 的最大值是 \blank{$\frac54+\frac{2\sqrt{3}}{9}$}．
	
	4．从 9 位同学中选出 5 位组成班委会，要求甲、乙二人或者同时入选，或者同时不入选，丙、丁二人不同时入选，则符合要求的选法共有 \blank{$41$} 种（用数字作答）。
	
	5．已知在等腰 $\triangle A B C$ 中，$A C=B C=\sqrt{5}$ ，点 $D, E, F$ 分别在边 $A B, B C, C A$ 上，且 $A D=D B=E F=1$ ．若 $\overrightarrow{D E} \cdot \overrightarrow{D F} \leqslant \frac{25}{16}$ ，则 $\overrightarrow{E F} \cdot \overrightarrow{B A}$ 的取值范围是 \blank{$\left[\frac43,2\right]$}．
	
	6．圆周上每个点都被染为红、黄、蓝三色之一，并且三种颜色的点都有。现从圆周上任取 $n$ 个点，若其中总存在三个点构成三个顶点同色的钝角三角形，则 $n$ 的最小可能值为
	\blank{13}．
	
	7．已知 $\alpha, \beta \in[0, \pi]$ ，则 $(\sin \alpha+\sin (\alpha+\beta)) \cdot \sin \beta$ 的最大值是 \blank{$\frac{8\sqrt{3}}{9}$}．
	
	8．设 $\frac{1}{1-x-x^2-x^3}=\sum_{n=0}^{\infty} a_n x^n$ ，若 $a_{n-1}=n^2$ ，则 $n$ 的取值集合为 \blank{$\{1,9\}$}．
	
	
	
\end{CJK}
\end{document}
